%%%%%%%%%%%%%%%%%%%%%%%%%%%%%%%%%%%%%%%%%
% Adaptive Time-Stepping Poster
% Based on Jacobs Landscape Poster template
%%%%%%%%%%%%%%%%%%%%%%%%%%%%%%%%%%%%%%%%%

\documentclass[final]{beamer}

\usepackage[scale=1.24]{beamerposter}
\usetheme{confposter}

\usepackage{graphicx}
\usepackage{booktabs}
\usepackage{amsmath, amssymb}
\usepackage{newtxmath}

\setbeamercolor{block title}{fg=black,bg=white}
\setbeamercolor{block body}{fg=black,bg=white}
\setbeamercolor{block alerted title}{fg=white,bg=black}
\setbeamercolor{block alerted body}{fg=black,bg=white}

\newlength{\sepwid}
\newlength{\onecolwid}
\newlength{\twocolwid}
\newlength{\threecolwid}
\setlength{\paperwidth}{48in}
\setlength{\paperheight}{38in}
\setlength{\sepwid}{0.024\paperwidth}
\setlength{\onecolwid}{0.22\paperwidth}
\setlength{\twocolwid}{0.464\paperwidth}
\setlength{\threecolwid}{0.708\paperwidth}
\setlength{\topmargin}{-0.5in}

\title{Adaptive Time-Stepping for Massively-Parallel Hardware-Accelerated Robot Simulation}
\author{Marco DiGiorgio and Vince Kurtz}
\institute{DePaul University}

\begin{document}

\addtobeamertemplate{block end}{}{\vspace*{2ex}}
\addtobeamertemplate{block alerted end}{}{\vspace*{2ex}}

\setlength{\belowcaptionskip}{2ex}
\setlength\belowdisplayshortskip{2ex}

% Thin border hugging each figure image.
\setlength{\fboxsep}{0pt}
\setlength{\fboxrule}{1.5pt}
\newcommand{\figimg}[1]{\fbox{\includegraphics[width=\dimexpr\linewidth-2\fboxrule\relax]{#1}}}

\begin{frame}[t]
\begin{columns}[t]

\begin{column}{\sepwid}\end{column}

%==========================================================
% LEFT COLUMN: motivation, hero scaling plot
%==========================================================
\begin{column}{\onecolwid}

\begin{block}{The Problem}
Massively parallel GPU simulators power modern robot learning, but every
world runs a \textbf{single fixed timestep}. Small $\delta t$ wastes work
everywhere to satisfy the stiffest contact; large $\delta t$ produces
penetration, jitter, and non-physical artifacts \textbf{that learned
policies exploit, breaking sim-to-real transfer}.

\medskip
\textbf{Error-controlled integration already solves this} by varying
$\delta t$ with system stiffness. But existing controllers are
\textbf{single-world and CPU-bound}: they assume one trajectory at a
time, not thousands of worlds stepping in lockstep on a GPU.
\end{block}

\begin{figure}
\centering
\figimg{SCALING.png}
\caption{\textbf{Wall-time scaling.} Per-world adaptive vs fixed-$\delta t$,
$N = 4$ to $8192$ worlds.}
\end{figure}

\begin{block}{Our Contribution}
We bring convex error-controlled integration\footnotesize$^{1}$\normalsize\
\textbf{from the CPU to massively parallel GPU hardware}. Each of $N$
worlds advances its \textbf{own inner $\delta t$} between fixed outer
control boundaries, with step acceptance, error estimation, and step-size
update all running on-device. The only host sync is a 4-byte boundary
flag per iteration; no per-state PCIe transfers in the physics loop.

\medskip
\footnotesize$^{1}$Convex error-controlled integration (CENIC),
\textit{arXiv:2511.08771}.\normalsize
\end{block}

\end{column}

\begin{column}{\sepwid}\end{column}

%==========================================================
% CENTER COLUMN (two-wide): method + results
%==========================================================
\begin{column}{\twocolwid}

%----- top row of center column: method, two sub-columns -----
\begin{columns}[t,totalwidth=\twocolwid]

\begin{column}{\onecolwid}\vspace{-.6in}

\begin{block}{Per-World Adaptive Time Stepping}
The system evolves continuously according to
\[
\Large
x(t_{n+1}) = x(t_n) + \int_{t_n}^{t_{n+1}} f(\tau, x(\tau))\,d\tau,
\]
which the integrator approximates over $\delta t = t_{n+1} - t_n$.
Each of the $N$ parallel worlds carries its own state $x^{(i)}$ and
its own inner step $\delta t^{(i)}$, advancing independently between
fixed outer control boundaries.
\end{block}

\end{column}

\begin{column}{\onecolwid}\vspace{-.6in}

\begin{block}{Step Doubling}
Let $\Phi_{\delta t}(x_n)$ denote one MuJoCo Warp convex step. Compare a
full step against two successive half-steps:
\begin{align}
\hat{x}_{n+1} &= \Phi_{\delta t}(x_n), \\
x_{n+1}\;\;\, &= \Phi_{\delta t/2}\bigl(\Phi_{\delta t/2}(x_n)\bigr).
\end{align}
The step-doubling difference yields a second-order error estimate
\begin{equation}
e_{n+1} = \lVert S\,(q_{n+1} - \hat{q}_{n+1}) \rVert_{\infty}
       \approx \bar{c}\,\delta t^{\,p}, \quad p = 2,
\end{equation}
a weighted, position-only inf-norm per world with $S = \operatorname{diag}(M)^{-1/2}$.
\end{block}

\end{column}
\end{columns}

%----- middle row: step-size update spans full width -----
\setbeamercolor{block alerted title}{fg=black,bg=dalgold}
\setbeamercolor{block alerted body}{fg=black,bg=white}

\begin{alertblock}{Step-Size Selection}
Accept if $e_{n+1} \le \varepsilon_{\mathrm{acc}}$ (or the controller
grows $\delta t$); otherwise reject and retry. Update the step from the
asymptotic error model:
\[
\delta t_{\mathrm{new}} = k_{\mathrm{safe}}\,\delta t \left(\frac{\varepsilon_{\mathrm{acc}}}{e_{n+1}}\right)^{\!1/p},
\qquad
\delta t = \min\!\left(\delta t_{\mathrm{new}},\;
k_{\mathrm{grow}}\,\delta t,\; t_{\mathrm{outer}} - t_n\right).
\]
Constants: $k_{\mathrm{safe}} = 0.9$, $k_{\mathrm{grow}} = 5.0$. The
boundary clamp prevents fast worlds from overshooting, while slow worlds
catch up. A symmetric deadband $[0.9, 1.2]\,\delta t$ and shrink cap
$k_{\mathrm{shrink}} = 0.1$ suppress step-size thrash.
\end{alertblock}

%----- bottom row: two result figures -----
\begin{columns}[t,totalwidth=\twocolwid]

\begin{column}{\onecolwid}

\begin{figure}
\centering
\figimg{FORK_FIXED.png}
\caption{\textbf{Fixed step, $\delta t = 10$ ms.} Fork penetrates rack.}
\end{figure}

\end{column}

\begin{column}{\onecolwid}

\begin{figure}
\centering
\figimg{FORK_ADAPTIVE.png}
\caption{\textbf{Adaptive step, $\varepsilon_{\mathrm{acc}} = 10^{-4}$.}
Fork rests on rack.}
\end{figure}

\end{column}
\end{columns}

\end{column}

\begin{column}{\sepwid}\end{column}

%==========================================================
% RIGHT COLUMN: adaptive fork, results, contact
%==========================================================
\begin{column}{\onecolwid}

\begin{figure}
\centering
\figimg{WORK_PRECISION.png}
\caption{\textbf{Work-precision.} Per-world $\delta t$ beats global
worst-case $\delta t$ across all tolerances.}
\end{figure}

\begin{figure}
\centering
\figimg{FRANKA_FARM.png}
\caption{\textbf{1024 worlds in parallel.} Franka arms stepped
simultaneously on one GPU, each tinted by its per-world adaptive
$\delta t$ (green = large step / free motion, red = small step /
stiff contact). Worlds diverge and each finds its own $\delta t$.}
\end{figure}

\begin{block}{Key Results}
\begin{itemize}
\item Integration error held inside user-specified tolerance
\item Contact artifacts eliminated at comparable wall-time budgets
\end{itemize}
\end{block}

\setbeamercolor{block title}{fg=black,bg=white}

\setbeamercolor{block alerted title}{fg=black,bg=dalblue}
\setbeamercolor{block alerted body}{fg=black,bg=white}

\end{column}

\end{columns}
\end{frame}

\end{document}